\documentclass[12pt]{article}

% ============================================================
% Compile: pdflatex (mặc định của Overleaf) hoặc xelatex đều được.
% Cần file logo-bk.png đặt cùng thư mục với file .tex.
% Hình EDA sinh bởi eda.py vào Assignment_2/output/eda/; hình kết quả sinh bởi analyze.py vào output/a2_2/;
% bảng số sinh bởi analyze.py vào report/generated/.
% ============================================================
\usepackage{iftex}
\ifPDFTeX
  \usepackage[utf8]{inputenc}
  \usepackage[T5]{fontenc}
  \usepackage{lmodern}
\else
  \usepackage{fontspec}
\fi
\usepackage[vietnamese]{babel}

\usepackage[a4paper,margin=2.5cm,headheight=2cm,headsep=0.8cm]{geometry}
\usepackage{graphicx,booktabs,tabularx,xltabular,longtable,array,xcolor,enumitem,amssymb,amsmath,multirow,caption}
\usepackage{fancyhdr}
\usepackage{lastpage}
\usepackage{titlesec}
\usepackage{float}
\usepackage{hyperref}
\hypersetup{hidelinks, pdftitle={A2.2 - Pipeline, ket qua va phan tich paper SegFormer - Phan doan 38 lop Oxford-IIIT Pet}}
\usepackage{parskip}

\graphicspath{{./}{figures/}}
\captionsetup{labelfont=bf,skip=4pt}
\setlist{nosep,leftmargin=*}
\titleformat{\section}{\large\bfseries}{\thesection.}{0.4em}{}
\titleformat{\subsection}{\normalsize\bfseries}{\thesubsection.}{0.4em}{}
\titlespacing*{\section}{0pt}{10pt}{4pt}\titlespacing*{\subsection}{0pt}{6pt}{2pt}
\newcolumntype{L}{>{\raggedright\arraybackslash}X}
\newcolumntype{C}{>{\centering\arraybackslash}X}
\newcolumntype{P}[1]{>{\raggedright\arraybackslash}p{#1}}
\newcolumntype{R}[1]{>{\raggedleft\arraybackslash}p{#1}}
\renewcommand{\arraystretch}{1.15}
\setlength{\emergencystretch}{3em}
\newcommand{\code}[1]{\texttt{#1}}
\newcommand{\st}[1]{\textbf{#1}}
\newcommand{\todo}[1]{\textcolor{red}{\textbf{[#1]}}}
% Chèn hình nếu có file, nếu chưa upload thì hiện khung báo thiếu (để vẫn compile được)
\newcommand{\fig}[2]{\IfFileExists{#2}{\includegraphics[width=#1]{#2}}{%
  \IfFileExists{../output/eda/#2}{\includegraphics[width=#1]{../output/eda/#2}}{%
  \IfFileExists{../output/a2_2/#2}{\includegraphics[width=#1]{../output/a2_2/#2}}{%
  \fbox{\parbox{0.8\linewidth}{\centering\todo{Thiếu hình \detokenize{#2}}}}}}}}

% ---------- Thông tin chung ----------
\newcommand{\TenTruong}{Đại học Quốc gia Thành phố Hồ Chí Minh\\Trường Đại học Bách khoa\\Khoa Khoa học và Kỹ thuật Máy tính}
\newcommand{\TenTruongHeader}{Ho Chi Minh University of Technology\\Faculty of Computer Science and Engineering}
\newcommand{\TenMonHoc}{Học sâu và ứng dụng trong thị giác máy tính}
\newcommand{\MaMonHoc}{CO5085}
\newcommand{\TenBaoCao}{Bài tập lớn A2 - Mốc A2.2: Pipeline, kết quả và phân tích paper}
\newcommand{\TenDeTai}{Phân đoạn ngữ nghĩa 38 lớp trên Oxford-IIIT Pet:\\So sánh U-Net (CNN) và SegFormer (Transformer)}
\newcommand{\GiangVien}{Thầy Lê Thành Sách\\Thầy Nguyễn Quốc Minh}
\newcommand{\ThanhVien}{%
  Đàm Quang Phục - 2570483\\
  Hồ Bảo An - 2570164
}
\newcommand{\DiaDiemNgay}{Thành Phố Hồ Chí Minh, 11/2026}
\newcommand{\FooterText}{\TenMonHoc\ -- \MaMonHoc}

% ---------- Header / Footer ----------
\pagestyle{fancy}
\fancyhf{}
\renewcommand{\headrulewidth}{0.4pt}
\renewcommand{\footrulewidth}{0.4pt}
\fancyhead[L]{%
  \parbox[t][\headheight][t]{\headwidth}{%
    \vspace{0.2cm}
    \begin{minipage}[b]{1.1cm}
      \IfFileExists{logo-bk.png}{\includegraphics[width=1.1cm]{logo-bk.png}}{}
    \end{minipage}%
    \hspace{0.5em}%
    \begin{minipage}[b]{8cm}
      \small \TenTruongHeader
    \end{minipage}%
  }%
}
\fancyfoot[L]{\small \FooterText}
\fancyfoot[R]{\small Page \thepage/\pageref{LastPage}}

\fancypagestyle{plain}{%
  \fancyhf{}
  \renewcommand{\headrulewidth}{0pt}
  \renewcommand{\footrulewidth}{0pt}
}

\begin{document}

% ================= TRANG BÌA =================
\begin{titlepage}
\thispagestyle{plain}
\begin{center}
{\bfseries \TenTruong \par}
\vspace{2em}

\IfFileExists{logo-bk.png}{\includegraphics[width=3cm]{logo-bk.png}\\[1em]}{}

\vspace{5em}

{\LARGE \bfseries \TenMonHoc\ (\MaMonHoc) \par}
\vspace{0.3em}
\rule{0.85\textwidth}{0.5pt}

\vspace{1em}
{\Large \bfseries \TenBaoCao \par}
\vspace{0.6em}
{\large \TenDeTai \par}
\vspace{0.3em}
\rule{0.85\textwidth}{0.5pt}

\vspace{5em}

\noindent\begin{minipage}[t]{3.6cm}
\raggedleft
Giảng viên:\\[0.5em]
\phantom{x}\\[0.5em]
Thành viên nhóm:
\end{minipage}%
\hspace{0.6em}%
\begin{minipage}[t]{8cm}
\GiangVien\\[0.5em]
\ThanhVien
\end{minipage}

\vspace{2em}
{\footnotesize Mã nguồn: \url{https://github.com/dqpinfo/Deep_Learning_Computer_Vision_CO5085/tree/main/Assignment_2}\par}

\vfill
{\bfseries \DiaDiemNgay}
\end{center}
\end{titlepage}

% ================= MỤC LỤC =================
\tableofcontents
\newpage

\newcommand{\gen}[1]{\IfFileExists{generated/#1}{\input{generated/#1}}{\todo{Thiếu generated/\detokenize{#1}} \\}}


% ======================================================================= tóm tắt
\noindent\fcolorbox{black!35}{black!3}{\parbox{\dimexpr\linewidth-2\fboxsep-2\fboxrule}{
\textbf{Tóm tắt mốc A2.2 (hạn 16/11/2026).} Nhóm huấn luyện U-Net + ResNet-50 và SegFormer-B2 với cùng công thức
(của paper SegFormer, 15 epoch, GPU T4) cho bài toán phân đoạn 38 lớp trên Oxford-IIIT Pet, ở 4 lượng dữ liệu train.
SegFormer-B2 đạt \st{89,2\% mIoU test} so với \st{46,1\%} của U-Net. Hai mô hình khoanh con vật tốt gần như nhau (IoU thú/nền
98,1\% và 96,9\%); khác biệt nằm ở việc \st{nhận đúng giống} (94,7\% so với 72,4\% ảnh). Chỉ với 10\% ảnh train, SegFormer vẫn
đạt 83,4\% mIoU, còn U-Net còn 4,8\%. Nhóm kiểm chứng thêm hai tuyên bố của paper: SegFormer giữ mIoU khi đổi độ phân giải test
(giảm tối đa 2,4 điểm) và bền với nhiễu, nhưng không bền với làm mờ hay nén JPEG nặng; ngược với paper, trên T4 nó chậm hơn và
tốn bộ nhớ hơn U-Net. Hạn chế chính: U-Net chưa hội tụ với công thức chung. Các bảng số được \code{analyze.py} sinh tự động từ log
JSON của từng lần chạy (\code{output/runs/}); số trong lời văn lấy từ các bảng đó.}}

% ======================================================================= 1
\section{Chủ đề và bài toán}
\subsection{Bài toán}
\st{Phân đoạn ngữ nghĩa} là gán cho \emph{mỗi điểm ảnh} một nhãn lớp. Khác với phân loại ảnh ở A1 (một nhãn cho cả ảnh),
mô hình phải trả lời đồng thời ``ở đâu'' và ``là gì''. Nhóm định nghĩa bài toán trên Oxford-IIIT Pet như sau:
\begin{itemize}
  \item \st{Đầu vào:} một ảnh RGB chứa một con chó hoặc mèo.
  \item \st{Đầu ra:} bản đồ nhãn cùng kích thước, mỗi điểm ảnh thuộc một trong 38 lớp: lớp 0 là nền, lớp 1--37 là giống
        của con vật (12 giống mèo, 25 giống chó).
  \item \st{Nhãn đúng:} ghép mặt nạ trimap chính thức (thú / nền / viền) với nhãn giống của ảnh. Điểm ảnh thuộc thú nhận
        mã giống, điểm nền nhận 0, vùng \emph{viền} (dải mờ quanh biên con vật mà người gán nhãn không chắc chắn) nhận
        giá trị 255 và \st{bị bỏ qua} khi tính loss và metric.
\end{itemize}

\subsection{Câu hỏi nghiên cứu}
\begin{enumerate}
  \item \st{CNN hay Transformer?} Với cùng nguồn pretrained (ImageNet-1k) và số tham số tương đương, U-Net + ResNet-50 hay
        SegFormer-B2 cho mIoU cao hơn trên bài toán 38 lớp, và với chi phí (thời gian, bộ nhớ, tốc độ) ra sao?
  \item \st{Lỗi đến từ đâu?} Phần mIoU bị mất là do vẽ sai biên hay do nhận sai giống? Mô hình nhầm giữa các giống cùng
        loài nhiều hơn nhầm chó $\leftrightarrow$ mèo không?
  \item \st{Mô hình nào cần ít nhãn hơn?} Khi chỉ còn 50\%, 25\%, 10\% ảnh train, khoảng cách giữa hai mô hình thay đổi thế nào?
  \item \st{Các tuyên bố của paper có đứng vững trên tập của nhóm không?} Ví dụ: bộ giải mã toàn MLP gọn nhẹ mà vẫn đủ tốt,
        encoder phân cấp không cần positional encoding (Mục~\ref{sec:paper}).
\end{enumerate}


% ======================================================================= 2
\section{Tập dữ liệu và EDA}
\subsection{Nguồn, license và cách tải}
Oxford-IIIT Pet được nhóm Visual Geometry Group (Đại học Oxford) và IIIT Hyderabad công bố cùng bài báo CVPR
2012~\cite{parkhi2012}, tại \url{https://www.robots.ox.ac.uk/~vgg/data/pets/}. Tác giả cung cấp dữ liệu
cho mục đích nghiên cứu; việc dùng ảnh phải tôn trọng điều khoản của các trang web gốc nơi ảnh được thu thập (theo file
\code{README} đi kèm bộ nhãn). Nhóm chỉ dùng cho học tập, không phân phối lại ảnh và trích dẫn bài báo gốc.

Mỗi ảnh có nhãn giống, nhãn loài (chó/mèo), hộp bao phần đầu và một mặt nạ \st{trimap} gán mỗi điểm ảnh là thú (1), nền (2)
hoặc viền chưa rõ (3). Hai danh sách chính thức chia dữ liệu thành \code{trainval} (3.680 ảnh) và \code{test} (3.669 ảnh).

Máy chủ của Oxford không truy cập được từ mạng của nhóm, nên \code{data.py} thử máy chủ gốc trước rồi tải hai file
\code{images.tar.gz} và \code{annotations.tar.gz} từ một bản sao trên Hugging Face. Hai file được \st{kiểm tra MD5} trùng
với giá trị thư viện torchvision dùng cho bản gốc, nên dữ liệu giống hệt bản chính thức. Thư mục \code{data/} không được
commit lên GitHub.

\subsection{Cách chia split}
\begin{table}[H]
\centering
\caption{Ba tập dữ liệu dùng chung cho mọi thí nghiệm (\code{splits/pet\_split.json}, \code{output/eda/summary.json}).}\label{tab:split}
\begin{tabularx}{\linewidth}{@{}lrrlL@{}}
\toprule
Split & Số ảnh & Mỗi giống & Nguồn & Vai trò \\\midrule
Train      & 3.303 & 80--90 & \code{trainval} chính thức & Cập nhật trọng số \\
Validation &   370 & 10     & \code{trainval} chính thức & Chọn checkpoint tốt nhất, theo dõi quá trình học \\
Test       & 3.662 & 88--100 & \code{test} chính thức    & Chỉ dùng một lần ở cuối mỗi thí nghiệm \\
\bottomrule
\end{tabularx}
\end{table}

Validation được lấy \st{phân tầng}: mỗi giống chọn ngẫu nhiên 10 ảnh từ danh sách \code{trainval} với seed 42
(\code{make\_split} trong \code{data.py}). Tên ảnh của cả ba split được lưu cố định trong \code{splits/pet\_split.json}.
Tập test chính thức giữ nguyên và chỉ dùng một lần ở cuối mỗi thí nghiệm. Thí nghiệm ít dữ liệu ở A2.2 lấy tập con của
train theo từng giống (\code{subsample}); tập nhỏ nằm trọn trong tập lớn hơn (10\% $\subset$ 25\% $\subset$ 50\%), validation
và test không đổi.

\subsection{Khám phá dữ liệu (EDA)}
Script \code{eda.py} quét toàn bộ 7.349 ảnh có trong hai danh sách chính thức (khoảng 20 giây trên Mac M1); kết quả lưu ở \code{output/eda/}.

\begin{figure}[H]
\centering
\fig{\linewidth}{breed_counts.png}
\caption{Số ảnh mỗi giống theo split. Tên in đậm là giống mèo.}\label{fig:counts}
\end{figure}

\textbf{Nhận xét về số lượng ảnh.}
\begin{itemize}
  \item Tập có 7.349 ảnh của \st{37 giống: 12 giống mèo và 25 giống chó}, mỗi giống khoảng 200 ảnh, chia gần đều giữa
        trainval và test (Hình~\ref{fig:counts}). Vì số giống chó gấp đôi giống mèo, số ảnh chó cũng gấp đôi
        (train: 1.065 mèo, 2.238 chó).
  \item Dữ liệu \st{cân bằng theo giống}: mỗi giống có 80--90 ảnh train, 10 ảnh validation, 88--100 ảnh test. Ít nhất là
        Egyptian Mau và Bombay, do tập gốc có ít ảnh hơn và do bước loại ảnh lỗi bên dưới.
  \item \st{Phát hiện lỗi nhãn:} khi quét trimap, nhóm thấy \st{14 ảnh} (7 trong trainval, 7 trong test) có trimap gán
        \emph{toàn bộ} ảnh là nền, dù ảnh rõ ràng có con vật (ví dụ \code{saint\_bernard\_108}, \code{Egyptian\_Mau\_162},
        \code{Persian\_259}). Giữ lại các ảnh này sẽ dạy mô hình rằng con vật là nền, và trừ điểm mô hình khi nó dự đoán đúng
        trên test. Nhóm \st{loại 14 ảnh này khỏi mọi split} và lưu danh sách trong \code{splits/pet\_split.json}
        (\code{excluded\_empty\_trimap}). Vì vậy split có 3.303 / 370 / 3.662 ảnh thay vì 3.310 / 370 / 3.669.
\end{itemize}

\begin{figure}[H]
\centering
\fig{0.95\linewidth}{samples_with_masks.png}
\caption{Ảnh mẫu (6 mèo, 6 chó trong tập train) và mặt nạ 38 lớp tương ứng.}\label{fig:samples}
\end{figure}

\begin{figure}[H]
\centering
\fig{\linewidth}{image_sizes.png}
\caption{Kích thước và tỉ lệ khung hình của ảnh gốc.}\label{fig:sizes}
\end{figure}

\begin{figure}[H]
\centering
\fig{\linewidth}{pixel_fractions.png}
\caption{Phân bố tỉ lệ điểm ảnh thú / nền / viền trong mỗi ảnh (đường đứt nét: trung vị).}\label{fig:fractions}
\end{figure}

\begin{figure}[H]
\centering
\fig{\linewidth}{class_pixel_share.png}
\caption{Tỉ lệ điểm ảnh có nhãn của mỗi lớp trên tập train (thang log), các giống xếp giảm dần.}\label{fig:share}
\end{figure}

\textbf{Nhận xét về ảnh và mặt nạ.}
\begin{itemize}
  \item \st{Kích thước ảnh không đồng nhất} (Hình~\ref{fig:sizes}): chiều rộng từ 114 đến 3.264 điểm ảnh, chiều cao từ 103
        đến 2.606, trung vị $500\times375$; khoảng 98\% ảnh có cạnh dài không quá 500. Tỉ lệ khung hình đa dạng, cả ngang lẫn dọc.
        Vì vậy cần đưa về một kích thước chung ($384\times384$) để ghép batch. Mọi trimap đều cùng kích thước với ảnh.
  \item 7.346 ảnh là RGB, 3 ảnh có kênh alpha (RGBA); \code{data.py} luôn chuyển về RGB.
  \item \st{Con vật thường chiếm phần nhỏ của ảnh} (Hình~\ref{fig:fractions}): trung vị 27\% điểm ảnh là thú, 60\% là nền và
        11\% là viền. Dải viền khá dày (Hình~\ref{fig:samples}) và bị bỏ qua, nên metric đo chất lượng ở phần thân con vật
        và nền rõ ràng, không đo độ sắc nét ở sát biên. Đây là giới hạn của nhãn gốc mà nhóm sẽ nêu khi thảo luận.
  \item \st{Mất cân bằng mạnh ở mức điểm ảnh} (Hình~\ref{fig:share}): trên tập train, nền chiếm 66\% điểm ảnh có nhãn, mỗi
        giống chỉ 0,4--1,5\%. Mô hình chỉ đoán ``nền'' ở mọi nơi đã đạt pixel accuracy 66\% nhưng mIoU chỉ khoảng
        $66/38\approx1,7\%$. Đây là lý do nhóm chọn mIoU (trung bình đều các lớp) làm metric chính thay vì pixel accuracy.
  \item Trung bình pixel trên 500 ảnh train ngẫu nhiên (R, G, B) là (0,465; 0,430; 0,382), độ lệch chuẩn (0,265; 0,260; 0,266),
        gần thống kê ImageNet, nên dùng chuẩn hoá ImageNet là hợp lý.
  \item Trong một ảnh chỉ có một con vật và một giống, nên mô hình cần nhận ra giống từ toàn cảnh con vật rồi tô nhất quán
        cả vùng. Một lỗi điển hình có thể đoán trước là một con vật bị tô bằng hai giống khác nhau ở hai phần cơ thể.
\end{itemize}

% ======================================================================= 3
\section{DataLoader và augmentation}\label{sec:loader}
\begin{table}[H]
\centering
\caption{Pipeline tiền xử lý, giống hệt nhau cho cả hai mô hình (\code{data.py}, torchvision \code{transforms.v2}).}\label{tab:pipeline}
\begin{tabularx}{\linewidth}{@{}P{3.3cm}LL@{}}
\toprule
Bước & Train & Validation / Test \\\midrule
Đọc dữ liệu & \multicolumn{2}{P{11.5cm}}{Ảnh JPEG chuyển sang RGB; trimap PNG chuyển thành mặt nạ 38 lớp (nền 0, giống 1--37, viền 255)} \\
Hình học (ảnh và mặt nạ cùng lúc) & \code{RandomResizedCrop(384, scale=(0.5, 1))} + \code{RandomHorizontalFlip} & \code{Resize} thẳng về $384\times384$ \\
Màu (chỉ ảnh) & \code{ColorJitter(0.2, 0.2, 0.2)} & Không \\
Chuẩn hoá & \multicolumn{2}{l}{Về $[0,1]$, rồi \code{Normalize} theo mean/std của ImageNet} \\
DataLoader & batch 16, \code{shuffle}, \code{drop\_last} & batch 16, không shuffle \\
\bottomrule
\end{tabularx}
\end{table}

\begin{figure}[H]
\centering
\fig{\linewidth}{augmentation_examples.png}
\caption{Cùng một ảnh: bản dùng cho validation/test (trái) và bốn bản augmentation cho train, kèm mặt nạ đi theo.}\label{fig:aug}
\end{figure}

\textbf{Lý do thiết kế.}
\begin{itemize}
  \item \st{Biến đổi hình học áp dụng đồng thời cho ảnh và mặt nạ}, và mặt nạ luôn được nội suy theo điểm gần nhất
        (nearest) để không sinh ra mã lớp không tồn tại ở biên. \code{transforms.v2} của torchvision xử lý việc này khi mặt
        nạ được bọc trong kiểu \code{tv\_tensors.Mask}; Hình~\ref{fig:aug} xác nhận mặt nạ khớp ảnh sau augmentation.
  \item \st{Kích thước $384\times384$:} ảnh gốc có cạnh khoảng 300--500 điểm ảnh (Mục EDA), nên 384 giữ được phần lớn chi
        tiết mà vẫn vừa bộ nhớ GPU T4 với batch 16. Cả hai mô hình nhận đúng cùng đầu vào; metric tính ở độ phân giải này.
  \item \st{Không lật dọc, không xoay mạnh:} ảnh thú cưng hầu như luôn đứng đúng chiều, nên các phép này chỉ tạo ảnh phi thực tế.
  \item \st{Chuẩn hoá ImageNet:} cả hai encoder được huấn luyện trước trên ImageNet-1k với cùng chuẩn hoá này.
\end{itemize}

% ======================================================================= 4
% Paper analysis section of the A2.2 report (SegFormer, Xie et al., NeurIPS 2021).
% Every number about the paper below was checked against the arXiv version (2105.15203).
\section{Phân tích paper SegFormer}\label{sec:paper}

\subsection{Thông tin và lý do chọn}
E.~Xie, W.~Wang, Z.~Yu, A.~Anandkumar, J.~M.~Alvarez, P.~Luo. \emph{SegFormer: Simple and Efficient Design for Semantic
Segmentation with Transformers}. NeurIPS 2021~\cite{xie2021}. Mã nguồn chính thức: \url{https://github.com/NVlabs/SegFormer};
trọng số encoder dùng trong bài: \code{nvidia/mit-b2} trên Hugging Face~\cite{wolf2020}. Nhóm chọn paper này vì một trong hai
mô hình của nhóm chính là SegFormer-B2, nên các nhận định của paper kiểm chứng được trực tiếp bằng thí nghiệm của nhóm.

\subsection{Vấn đề paper giải quyết}
Trước SegFormer, các mô hình phân đoạn dùng Transformer gặp hai trở ngại. Thứ nhất, ViT gốc chỉ cho đặc trưng ở một tỉ lệ
và độ phân giải thấp, trong khi phân đoạn cần cả chi tiết nhỏ (biên vật) lẫn ngữ cảnh rộng. Thứ hai, ViT dùng
positional encoding có kích thước cố định; khi ảnh test lớn hơn ảnh train, positional encoding phải nội suy và độ chính
xác giảm. Các hướng đi trước đó (SETR, PVT, Swin) chủ yếu thiết kế lại encoder và vẫn dùng bộ giải mã nặng. SegFormer
đặt giả thiết: \st{nếu encoder đủ tốt ở nhiều tỉ lệ thì bộ giải mã có thể rất nhẹ}.

\subsection{Phương pháp}
\textbf{Encoder Mix Transformer (MiT).} Gồm 4 tầng cho ra đặc trưng ở 1/4, 1/8, 1/16, 1/32 độ phân giải ảnh, giống kim
tự tháp đặc trưng của CNN. Mỗi tầng có ba điểm khác ViT:
\begin{itemize}
  \item \st{Overlapped patch merging:} cắt patch bằng tích chập có bước trượt nhỏ hơn kích thước nhân (ví dụ nhân 7, bước 4
        ở tầng đầu), nên các patch chồng lên nhau và giữ được tính liên tục cục bộ quanh biên patch.
  \item \st{Efficient self-attention:} trước khi tính attention, chuỗi key/value được rút ngắn $R$ lần ($R = 64, 16, 4, 1$
        ở 4 tầng, tức giảm 8, 4, 2, 1 lần theo mỗi chiều không gian), giảm chi phí từ $O(N^2)$ xuống $O(N^2/R)$. Nhờ vậy attention toàn cục vẫn
        chạy được ở độ phân giải 1/4.
  \item \st{Mix-FFN:} chèn một tích chập depthwise $3\times3$ vào khối FFN. Tích chập có zero-padding ``rò rỉ'' thông tin
        vị trí, nên mô hình \st{không cần positional encoding} và không bị ràng buộc vào kích thước ảnh lúc train.
\end{itemize}
\textbf{Bộ giải mã toàn MLP.} Bốn bản đồ đặc trưng được chiếu về cùng $C$ kênh bằng một lớp tuyến tính, phóng lên 1/4 độ
phân giải, ghép lại, trộn bằng một lớp tuyến tính nữa rồi dự đoán $N_{cls}$ lớp. Với B2, $C = 768$ và bộ giải mã chỉ có
3,3 triệu tham số. Paper lập luận rằng tầng nông của MiT đã nhìn cục bộ còn tầng sâu nhìn gần như toàn ảnh (minh hoạ
bằng effective receptive field), nên ghép đặc trưng các tầng là đủ, không cần module ngữ cảnh như ASPP.

\textbf{Loss và huấn luyện.} Chỉ dùng cross-entropy theo điểm ảnh; paper nói rõ không dùng OHEM, loss phụ hay cân bằng lớp.
AdamW, learning rate $6\times10^{-5}$, lịch ``poly'' hệ số 1,0; encoder pretrained ImageNet-1k, bộ giải mã khởi tạo ngẫu
nhiên. Augmentation: thay đổi kích thước ngẫu nhiên tỉ lệ 0,5--2,0, lật ngang, cắt ngẫu nhiên $512\times512$ (ADE20K,
COCO-Stuff) hoặc $1024\times1024$ (Cityscapes); 160 nghìn vòng lặp, batch 16, trên 8 GPU V100.

\subsection{Thí nghiệm của paper}
Ba tập: ADE20K (150 lớp, 20.210 ảnh), Cityscapes (19 lớp, 5.000 ảnh độ phân giải cao) và COCO-Stuff (172 lớp, 164 nghìn
ảnh); metric mIoU, kèm số tham số, FLOPs và tốc độ. Các kết quả nhóm thấy quan trọng nhất (Bảng~\ref{tab:paper-numbers}):

\begin{table}[H]
\centering
\caption{Một số kết quả của paper được nhóm đối chiếu (số lấy từ Bảng 1 và Bảng 5 của paper).}\label{tab:paper-numbers}
\begin{tabularx}{\linewidth}{@{}P{4.6cm}L@{}}
\toprule
Thí nghiệm & Kết quả \\\midrule
Cỡ mô hình (Bảng 1a) & SegFormer-B2: encoder 24,2M + decoder 3,3M tham số; mIoU 46,5 (ADE20K), 81,0 (Cityscapes), 44,6 (COCO-Stuff), test một tỉ lệ \\
Bộ giải mã MLP với encoder CNN (Bảng 1d, ADE20K) & ResNet-50 + bộ giải mã MLP: 34,7 mIoU; MiT-B2 + cùng bộ giải mã: 45,4. Bộ giải mã nhẹ chỉ hiệu quả khi encoder là Transformer \\
Mix-FFN và positional encoding (Bảng 1c, Cityscapes) & Train $768\times768$, test $1024\times2048$: dùng positional encoding giảm 77,3 $\rightarrow$ 74,0; dùng Mix-FFN chỉ giảm 80,5 $\rightarrow$ 79,8 \\
Số kênh $C$ của bộ giải mã (Bảng 1b) & 256 $\rightarrow$ 44,9; 768 $\rightarrow$ 45,4; 2048 $\rightarrow$ 45,6: tăng kênh ít tác dụng \\
Độ bền với ảnh hỏng (Bảng 5, Cityscapes-C) & SegFormer-B5 so với DeepLabv3+ (ResNet-50): ảnh gốc 82,4 / 76,6; nhiễu Gauss 57,8 / 6,5; làm mờ Gauss 69,8 / 57,7; nén JPEG 68,4 / 42,0 \\
\bottomrule
\end{tabularx}
\end{table}

\textbf{Kết luận của paper.} SegFormer đạt mIoU cao hơn các phương pháp cùng thời với ít tham số và FLOPs hơn; encoder không
positional encoding chịu được việc đổi độ phân giải lúc test; và mô hình bền hơn hẳn CNN khi ảnh bị nhiễu.

\subsection{Hạn chế của paper}
\begin{itemize}
  \item \st{So sánh độ bền không cùng cỡ.} Bảng 5 chỉ thử SegFormer-B5 (khoảng 85 triệu tham số) với các mô hình nhỏ hơn
        nhiều; paper không cho biết B2 có còn bền như vậy không, cũng không so với một CNN cùng cỡ được huấn luyện cùng
        công thức. Thí nghiệm ở Mục~\ref{sec:robust} của nhóm đo đúng điểm này.
  \item \st{Công thức huấn luyện khác nhau giữa các mô hình so sánh.} Kết quả của mô hình khác lấy lại từ paper gốc, với
        lịch huấn luyện và augmentation khác, nên chênh lệch không chỉ do kiến trúc.
  \item \st{Phụ thuộc dữ liệu và tính toán lớn.} Encoder cần pretrain ImageNet; mọi thí nghiệm dùng tập lớn (vài nghìn đến
        hàng trăm nghìn ảnh) và 8 GPU. Paper không thử chế độ ít nhãn.
  \item \st{Biên vật.} Bộ giải mã dự đoán ở 1/4 độ phân giải rồi phóng bilinear; với vật nhỏ hoặc biên mảnh (lông, tai) điều
        này có thể làm mất chi tiết. Paper chỉ minh hoạ định tính trên Cityscapes.
  \item Tác giả tự nêu: chưa rõ mô hình nhỏ nhất (3,7M tham số) có chạy được trên chip biên chỉ có 100~KB bộ nhớ không.
\end{itemize}

\subsection{Điều kiện trên tập của nhóm khác với paper}
Oxford-IIIT Pet khác các tập của paper ở ba điểm: (1) mỗi ảnh chỉ có \st{một con vật và nền}, nên ngữ cảnh toàn cảnh mà
attention mạnh ở ADE20K ít quan trọng hơn; (2) phần khó là \st{phân biệt 37 giống rất giống nhau}, gần với phân loại tinh
hơn là phân đoạn cảnh; (3) tập train chỉ có 3.303 ảnh, nhỏ hơn ADE20K khoảng 6 lần. Nhóm cũng huấn luyện ngắn hơn nhiều
(15 epoch, khoảng 3 nghìn vòng lặp) ở $384\times384$ trên một GPU T4.


% ======================================================================= 5
\section{Pipeline, huấn luyện và cách so sánh}\label{sec:pipeline}
\subsection{Hai mô hình}
\begin{table}[H]
\centering
\caption{Hai mô hình so sánh; số tham số đo bằng \code{smoke\_test.py} với đầu ra 38 lớp (\code{output/protocol\_table.json}).}\label{tab:models}
\begin{tabularx}{\linewidth}{@{}P{2.4cm}P{4.4cm}R{2.3cm}L@{}}
\toprule
Họ & Mô hình & Tham số & Đặc điểm \\\midrule
CNN & U-Net + ResNet-50 \newline \code{resnet50.a1\_in1k} (timm) & 32.521.606 & Encoder 23.508.032 + bộ giải mã 9.013.574; giải mã dần tới độ phân giải gốc, có skip connection \\
Transformer & SegFormer-B2 \newline \code{nvidia/mit-b2} (Hugging Face) & 27.375.846 & Encoder 24.196.288 + bộ giải mã MLP 3.179.558; dự đoán ở 1/4 độ phân giải \\
\bottomrule
\end{tabularx}
\end{table}

\begin{itemize}
  \item \st{U-Net + ResNet-50}~\cite{ronneberger2015,he2016}: encoder ResNet-50 (\code{timm}~\cite{timm}, trọng số
        \code{resnet50.a1\_in1k}) cho 5 bản đồ đặc trưng ở 1/2 đến 1/32 độ phân giải. Bộ giải mã do nhóm tự viết
        (\code{models.py}): 5 bước, mỗi bước phóng gấp đôi, ghép với đặc trưng encoder cùng tỉ lệ (skip connection) và qua hai
        lớp tích chập $3\times3$ + BatchNorm + ReLU, cuối cùng là lớp $1\times1$ ra 38 kênh ở độ phân giải gốc.
  \item \st{SegFormer-B2}~\cite{xie2021}: encoder MiT-B2 pretrained ImageNet-1k (\code{nvidia/mit-b2}) và bộ giải mã toàn MLP,
        lấy từ thư viện Hugging Face Transformers~\cite{wolf2020}. Đầu ra ở 1/4 độ phân giải được phóng bilinear về kích
        thước ảnh trước khi tính loss.
\end{itemize}
Hai mô hình có \st{cùng nguồn pretrained} (chỉ ImageNet-1k) và \st{số tham số cùng bậc}, để khác biệt kết quả đến từ kiến trúc
chứ không phải từ dữ liệu pretrain hay kích thước. Nhóm cố ý không dùng các trọng số đã được huấn luyện cho phân đoạn trên
COCO hay Pascal VOC, vì các tập này có sẵn lớp chó và mèo.

\subsection{Metric}\label{sec:metrics}
Mọi metric được tính trên \st{toàn bộ split} từ một ma trận nhầm lẫn $38\times38$ cộng dồn qua các ảnh (\code{metrics.py}),
bỏ qua điểm ảnh viền. Gọi $TP_c$, $FP_c$, $FN_c$ là số điểm ảnh đúng, dự đoán thừa và bỏ sót của lớp $c$:
\begin{align*}
\mathrm{IoU}_c &= \frac{TP_c}{TP_c + FP_c + FN_c}, &
\mathrm{mIoU} &= \frac{1}{38}\sum_{c=0}^{37}\mathrm{IoU}_c, &
\mathrm{Dice}_c &= \frac{2\,TP_c}{2\,TP_c + FP_c + FN_c}.
\end{align*}
\begin{itemize}
  \item \st{mIoU 38 lớp (metric chính):} mỗi lớp có trọng số như nhau, nên giống hiếm hay khó cũng ảnh hưởng ngang lớp nền.
        Đây là metric chuẩn của bài toán phân đoạn ngữ nghĩa và cũng là metric chính của paper SegFormer.
  \item \st{Mean Dice:} trung bình $\mathrm{Dice}_c$, metric thông dụng thứ hai mà đề bài nêu.
  \item \st{Pixel accuracy:} tỉ lệ điểm ảnh đúng. Chỉ báo cáo để tham khảo, vì nền chiếm phần lớn điểm ảnh (Hình~\ref{fig:share})
        nên metric này dễ cao dù nhận sai giống.
  \item \st{IoU thú/nền:} gộp 37 giống thành một lớp ``thú'' rồi tính IoU. So với mIoU, nó tách riêng chất lượng khoanh hình
        dạng khỏi chất lượng nhận giống.
  \item \st{Chi phí:} số tham số, thời gian huấn luyện, bộ nhớ GPU đỉnh và tốc độ suy luận (ảnh/giây) trên cùng GPU T4.
  \item \st{Độ chính xác giống theo ảnh:} mỗi ảnh test lấy giống được dự đoán nhiều nhất trong các điểm ảnh mô hình cho là
        thú, rồi so với giống thật. Metric này cho biết mô hình có ``hiểu'' đúng con vật không, tách khỏi chất lượng biên.
\end{itemize}

\subsection{Huấn luyện}\label{sec:train}
\begin{table}[H]
\centering
\caption{Siêu tham số, giống nhau cho cả hai mô hình và mọi lượng dữ liệu (\code{config.TRAIN\_HPARAMS}, \code{train.py}).}\label{tab:hparams}
\begin{tabularx}{\linewidth}{@{}P{3.6cm}L@{}}
\toprule
Thiết lập & Giá trị \\\midrule
Loss & Cross-entropy từng điểm ảnh, bỏ qua viền (\code{ignore\_index=255}), giống paper (không OHEM, không loss phụ) \\
Optimizer & AdamW, weight decay 0,01 \\
Learning rate & Encoder $6\times10^{-5}$ (giá trị của paper); bộ giải mã và lớp ra $6\times10^{-4}$ \\
Lịch learning rate & Warmup tuyến tính 1 epoch, sau đó ``poly'' luỹ thừa 1,0 giảm về 0 (giống paper) \\
Số epoch / batch / seed & 15 / 16 / 42; với 100\% dữ liệu là 206 bước mỗi epoch, khoảng 3.100 bước \\
Độ chính xác số & Mixed precision (fp16) trên một GPU Tesla T4 (Google Colab) \\
Chọn checkpoint & Epoch có mIoU validation cao nhất; sau đó đánh giá trên test đúng một lần \\
\bottomrule
\end{tabularx}
\end{table}

Kế hoạch ở A2.1 là 40 epoch. Khi đo trên T4, một epoch U-Net với 100\% dữ liệu mất khoảng 84 giây, nên 8 lần chạy
40 epoch cần hơn 4 giờ GPU, vượt hạn mức Colab miễn phí còn lại. Nhóm giảm xuống \st{15 epoch cho mọi cấu hình}, giữ nguyên
các thiết lập khác để so sánh vẫn công bằng. Lịch poly co theo số epoch nên learning rate vẫn giảm về 0 ở cuối.

Vòng lặp huấn luyện tự viết bằng PyTorch (\code{train.py}, dựa trên khung của A1): checkpoint sau mỗi epoch để chạy tiếp
khi Colab bị ngắt, log JSON theo epoch, đo thời gian, bộ nhớ đỉnh và tốc độ suy luận (20 batch sau khi khởi động).
\code{run\_all.py} chạy lần lượt cả 8 cấu hình và bỏ qua cấu hình đã xong; \code{colab\_a2\_2.ipynb} gói toàn bộ quy trình.

\subsection{Ma trận thí nghiệm}
\begin{table}[H]
\centering
\caption{8 lần huấn luyện và đánh giá thêm (\code{config.EXPERIMENTS}, \code{eval\_extra.py}).}\label{tab:matrix}
\begin{tabularx}{\linewidth}{@{}P{3.6cm}P{3.2cm}L@{}}
\toprule
Cấu hình & Ảnh train & Mục đích \\\midrule
\code{\{unet\_r50, segformer\_b2\}\_f100} & 3.303 (100\%) & So sánh chính CNN và Transformer \\
\code{\_f50}, \code{\_f25}, \code{\_f10} & 1.650 / 812 / 332 & Hiệu quả dữ liệu: mô hình nào giữ mIoU tốt hơn khi ít nhãn. Tập nhỏ nằm trong tập lớn, lấy đều theo giống \\
Đánh giá thêm (không huấn luyện lại) & mô hình 100\% & (1) Test ở $256$--$640$ điểm ảnh (train ở 384): kiểm chứng tuyên bố Mix-FFN không cần positional encoding;
(2) ảnh test bị nhiễu Gauss, làm mờ, nén JPEG, mỗi loại 3 mức: kiểm chứng tuyên bố về độ bền (Bảng 5 của paper) ở cỡ B2 \\
\bottomrule
\end{tabularx}
\end{table}

% ======================================================================= 6
\section{Kết quả}\label{sec:results}
\subsection{So sánh chính (100\% dữ liệu train)}
\begin{table}[H]
\centering
\caption{Kết quả trên tập test 3.662 ảnh (mIoU validation để chọn checkpoint). Đơn vị \%.}\label{tab:main}
\begin{tabularx}{\linewidth}{@{}LrCCCCCC@{}}
\toprule
Mô hình & Tham số & mIoU val & mIoU test & mDice & Pixel acc. & IoU thú/nền & Giống đúng/ảnh \\\midrule
% from generated/main_table.tex (analyze.py)
U-Net + ResNet-50 & 32,5M & 48,5 & 46,1 & 58,3 & 85,9 & 96,9 & 72,4 \\
SegFormer-B2 & 27,4M & 89,5 & 89,2 & 93,9 & 97,5 & 98,1 & 94,7 \\
\bottomrule
\end{tabularx}
\end{table}

\begin{table}[H]
\centering
\caption{Chi phí trên một GPU Tesla T4 (fp16).}\label{tab:cost}
\begin{tabularx}{\linewidth}{@{}LCCCCC@{}}
\toprule
Mô hình & Epoch tốt nhất & Giây/epoch & Tổng (phút) & Bộ nhớ đỉnh (GB) & Suy luận (ảnh/s) \\\midrule
% from generated/cost_table.tex (analyze.py)
U-Net + ResNet-50 & 15/15 & 93 & 23 & 4,6 & 134 \\
SegFormer-B2 & 11/15 & 128 & 32 & 8,3 & 77 \\
\bottomrule
\end{tabularx}
\end{table}

\begin{figure}[H]
\centering
\fig{\linewidth}{learning_curves.png}
\caption{Loss và mIoU validation theo epoch (100\% dữ liệu train).}\label{fig:curves}
\end{figure}

\subsection{Lỗi đến từ đâu}
\begin{figure}[H]
\centering
\fig{\linewidth}{species_confusion.png}
\caption{Ma trận nhầm lẫn gộp nền / mèo / chó trên tập test (chuẩn hoá theo hàng).}\label{fig:species}
\end{figure}

\begin{figure}[H]
\centering
\fig{\linewidth}{per_breed_iou.png}
\caption{IoU từng giống trên tập test.}\label{fig:breed}
\end{figure}

\begin{figure}[H]
\centering
\fig{\linewidth}{examples.png}
\caption{Ảnh test điển hình và bản đồ lỗi của hai mô hình.}\label{fig:examples}
\end{figure}

\subsection{Hiệu quả dữ liệu}
\begin{table}[H]
\centering
\caption{Kết quả test khi giảm dữ liệu train (validation và test giữ nguyên). Đơn vị \%.}\label{tab:fraction}
\begin{tabularx}{\linewidth}{@{}LCCCCCC@{}}
\toprule
& \multicolumn{3}{c}{U-Net + ResNet-50} & \multicolumn{3}{c}{SegFormer-B2} \\\cmidrule(lr){2-4}\cmidrule(lr){5-7}
Ảnh train & mIoU & IoU thú & Giống/ảnh & mIoU & IoU thú & Giống/ảnh \\\midrule
% from generated/fraction_table.tex (analyze.py)
100\% (3.303) & 46,1 & 96,9 & 72,4 & 89,2 & 98,1 & 94,7 \\
50\% (1.650) & 25,5 & 96,6 & 48,4 & 88,2 & 98,1 & 94,0 \\
25\% (812) & 12,4 & 94,7 & 28,2 & 85,0 & 97,7 & 92,5 \\
10\% (332) & 4,8 & 90,7 & 10,6 & 83,4 & 97,1 & 91,6 \\
\bottomrule
\end{tabularx}
\end{table}

\begin{figure}[H]
\centering
\fig{\linewidth}{data_efficiency.png}
\caption{Kết quả test theo số ảnh train.}\label{fig:efficiency}
\end{figure}

\subsection{Kiểm chứng tuyên bố của paper: độ phân giải và ảnh hỏng}\label{sec:robust}
\begin{table}[H]
\centering
\caption{mIoU test (\%) khi đổi kích thước ảnh test; cả hai mô hình huấn luyện ở $384\times384$.}\label{tab:resolution}
\begin{tabularx}{\linewidth}{@{}LCCCCC@{}}
\toprule
Mô hình & 256 & 320 & 384 & 512 & 640 \\\midrule
% from generated/resolution_table.tex (analyze.py)
U-Net + ResNet-50 & 48,5 & 46,3 & 43,4 & 37,5 & 32,3 \\
SegFormer-B2 & 86,9 & 88,2 & 89,0 & 88,4 & 86,6 \\
\bottomrule
\end{tabularx}
\end{table}

\begin{table}[H]
\centering\small
\caption{mIoU test (\%) trên ảnh bị hỏng, ba mức độ mỗi loại.}\label{tab:corruption}
\begin{tabularx}{\linewidth}{@{}LCCCCCCCCCC@{}}
\toprule
& & \multicolumn{3}{c}{Nhiễu Gauss $\sigma$} & \multicolumn{3}{c}{Làm mờ Gauss $\sigma$} & \multicolumn{3}{c}{JPEG chất lượng} \\
\cmidrule(lr){3-5}\cmidrule(lr){6-8}\cmidrule(lr){9-11}
Mô hình & Gốc & 0,04 & 0,08 & 0,16 & 1 & 2 & 4 & 30 & 15 & 5 \\\midrule
% from generated/corruption_table.tex (analyze.py)
U-Net + ResNet-50 & 43,4 & 36,0 & 21,8 & 2,6 & 38,6 & 24,7 & 12,2 & 31,2 & 19,4 & 3,6 \\
SegFormer-B2 & 89,0 & 88,9 & 87,5 & 79,5 & 87,4 & 77,8 & 47,2 & 86,8 & 79,5 & 45,6 \\
\bottomrule
\end{tabularx}
\end{table}

\begin{figure}[H]
\centering
\fig{\linewidth}{robustness.png}
\caption{mIoU theo độ phân giải test và mức độ hỏng ảnh.}\label{fig:robust}
\end{figure}

\begin{figure}[H]
\centering
\fig{\linewidth}{corruption_examples.png}
\caption{Một ảnh test dưới các loại và mức độ hỏng đã dùng.}\label{fig:corrupt-ex}
\end{figure}

% ======================================================================= 7
\section{Thảo luận và đối chiếu với paper}\label{sec:discussion}
\subsection{CNN hay Transformer?}
Với cùng dữ liệu, cùng công thức huấn luyện và 15 epoch, SegFormer-B2 đạt mIoU test 89,2\% còn U-Net + ResNet-50 chỉ
46,1\% (Bảng~\ref{tab:main}). Khoảng cách gần như nằm hết ở việc \st{nhận giống}, không phải ở việc khoanh con vật:
\begin{itemize}
  \item IoU thú/nền của hai mô hình gần nhau (98,1\% và 97,0\%). Cả hai đều tách con vật khỏi nền tốt; U-Net chỉ kém
        khoảng 1 điểm.
  \item Độ chính xác giống theo ảnh chênh lớn: 94,7\% so với 72,4\%. Khi U-Net đoán sai giống, cả vùng con vật bị tính
        sai lớp, nên IoU của lớp đúng và của lớp bị đoán nhầm cùng giảm. Đó là lý do mIoU chênh hơn 40 điểm dù hình dạng
        gần như đúng.
\end{itemize}
Đường học (Hình~\ref{fig:curves}) cho thấy hai mô hình học theo hai nhịp rất khác nhau. SegFormer đạt mIoU validation
80\% ngay sau epoch 1 và gần bão hoà từ epoch 3. U-Net chỉ đạt 3,5\% sau epoch 1, rồi tăng chậm và đều tới 48,5\% ở
epoch 15, khi đường cong vẫn còn dốc. Vì vậy \st{kết quả của U-Net là kết quả chưa hội tụ}: với nhiều epoch hơn hoặc
learning rate lớn hơn cho encoder ResNet, khoảng cách chắc chắn sẽ hẹp lại. Nhóm giải thích sự chênh lệch như sau:
\begin{itemize}
  \item \st{Công thức huấn luyện là của SegFormer.} Learning rate $6\times10^{-5}$ cho encoder vốn được paper chọn cho MiT.
        Với ResNet-50 (có BatchNorm, quen với learning rate lớn hơn nhiều), mức này làm encoder gần như đứng yên, nên
        U-Net phải dựa vào bộ giải mã 9 triệu tham số học từ đầu. Đây đúng là hạn chế nhóm đã nêu ở paper
        (Mục~\ref{sec:paper}): so sánh dùng cùng một công thức thì công thức đó có thể thiên vị một kiến trúc.
  \item \st{Nhận giống cần ngữ cảnh toàn cục.} Phân biệt 37 giống dựa vào toàn bộ con vật (mặt, lông, dáng). Self-attention
        của MiT nhìn toàn ảnh ngay từ tầng đầu, còn CNN chỉ có ngữ cảnh rộng ở tầng sâu. Trong U-Net, quyết định ở mỗi điểm
        ảnh còn chịu ảnh hưởng mạnh của skip connection từ tầng nông, nên một con vật có thể bị tô bằng nhiều giống
        (Hình~\ref{fig:examples}).
  \item \st{Bộ giải mã nhẹ vẫn đủ.} Bộ giải mã MLP của SegFormer chỉ 3,2 triệu tham số và dự đoán ở 1/4 độ phân giải, nhưng
        IoU thú/nền vẫn cao hơn U-Net. Trên tập này, nhãn viền bị bỏ qua nên lợi thế ``biên sắc nét'' của U-Net không được
        tính (Mục EDA). Điều này khớp với tuyên bố trung tâm của paper: encoder tốt thì bộ giải mã có thể rất đơn giản.
\end{itemize}

\subsection{Chi phí}
SegFormer-B2 ít tham số hơn (27,4 triệu so với 32,5 triệu) nhưng \st{đắt hơn khi chạy} trên T4 (Bảng~\ref{tab:cost}): mỗi
epoch 128 giây so với 93 giây, bộ nhớ đỉnh 8,3 GB so với 4,6 GB, suy luận 77 so với 134 ảnh/giây. Self-attention ở độ phân
giải 1/4 (dù đã rút gọn key/value) và đặc trưng 768 kênh ở 1/4 độ phân giải của bộ giải mã tốn bộ nhớ hơn tích chập.
Khi đánh giá ở $640\times640$, SegFormer với batch 16 hết bộ nhớ T4 ở lớp trộn của bộ giải mã (cần thêm 2,3 GB), nhóm phải giảm
batch xuống 4. Số tham số nhỏ của paper vì vậy không đồng nghĩa với rẻ hơn trên GPU phổ thông. Tuy vậy, tính theo
\st{chất lượng trên cùng thời gian}, SegFormer vẫn lợi hơn hẳn: sau 1 epoch (khoảng 2 phút) nó đã vượt xa U-Net sau 15 epoch.

\subsection{Lỗi đến từ đâu}
Ma trận gộp nền / mèo / chó (Hình~\ref{fig:species}) và phân tích điểm ảnh thú trên test cho thấy:
\begin{itemize}
  \item U-Net gán đúng giống cho 60,2\% điểm ảnh thú; 30,0\% bị gán \st{sai giống cùng loài}, 8,1\% bị gán \st{sai loài}
        (chủ yếu mèo bị đoán là chó: 15,0\% điểm ảnh mèo), chỉ 1,7\% bị đoán là nền. SegFormer gán đúng 93,8\%; sai
        giống cùng loài 5,4\%, sai loài chỉ 0,1\%, đoán là nền 0,7\%.
  \item Vậy lỗi của cả hai mô hình chủ yếu là \st{nhầm giữa các giống cùng loài}, không phải nhầm chó với mèo hay bỏ sót
        con vật, đúng như dự đoán ở A2.1. U-Net còn có lỗi đặc trưng là tô một con vật bằng nhiều giống (ảnh Bengal\_24 ở
        Hình~\ref{fig:examples}), điều gần như không xảy ra với SegFormer.
  \item Theo từng giống (Hình~\ref{fig:breed}), SegFormer tốt hơn U-Net ở \st{cả 37 giống}; chênh lệch nhỏ nhất vẫn là 18,7 điểm
        (newfoundland). Giống khó nhất với SegFormer là american pit bull terrier (IoU 58,9\%), Ragdoll (62,6\%) và
        staffordshire bull terrier (65,2\%): pit bull và staffordshire gần như giống hệt nhau, còn Ragdoll dễ lẫn với các
        giống mèo lông dài khác (Ragdoll\_62 bị đoán là Maine Coon). Ảnh american\_bulldog\_61 và Ragdoll\_62 là hai trường hợp U-Net đúng còn SegFormer nhầm sang
        giống họ hàng.
  \item \st{Nhãn sai trong tập test.} Nhóm chọn nhóm ảnh ``cả hai cùng sai'' để xem và thấy hai ảnh (Persian\_273,
        Maine\_Coon\_268) có trimap bị \st{đảo ngược}: con vật được gán là nền và nền được gán là thú. Cả hai mô hình đều
        khoanh đúng con vật nhưng bị chấm sai toàn bộ. Đây là loại lỗi nhãn thứ hai, khác với 14 ảnh ``toàn nền'' đã loại ở
        A2.1; nhóm không loại các ảnh này khỏi test sau khi đã thấy kết quả để tránh chỉnh tập test theo mô hình, nhưng
        ghi nhận rằng mIoU báo cáo hơi thấp hơn giá trị thật.
\end{itemize}

% ======================================================================= 8
\subsection{Mô hình nào cần ít nhãn hơn?}
Bảng~\ref{tab:fraction} và Hình~\ref{fig:efficiency} trả lời rõ: \st{SegFormer gần như không phụ thuộc vào lượng nhãn} trong
khoảng thử, còn U-Net sụp nhanh. Khi giảm từ 3.303 xuống 332 ảnh train (10\%), mIoU của SegFormer chỉ giảm từ 89,2\% xuống
83,4\% và độ chính xác giống theo ảnh từ 94,7\% xuống 91,6\%; với chỉ 9 ảnh mỗi giống, nó vẫn tốt hơn U-Net dùng toàn bộ dữ liệu.
U-Net giảm từ 46,1\% xuống 4,8\% mIoU (giống đúng theo ảnh 10,6\%, gần mức đoán bừa trong 37 giống là 2,7\%). Trong khi đó IoU
thú/nền của U-Net vẫn giữ trên 90\% ở mọi mức: tách con vật khỏi nền là việc dễ, cần ít nhãn; còn nhận giống mới là phần cần
dữ liệu, và ở đây encoder MiT pretrained cung cấp đặc trưng ``sẵn sàng'' cho phân biệt giống tốt hơn nhiều. Lưu ý mọi cấu hình
đều chạy 15 epoch, nên cấu hình ít ảnh có ít bước cập nhật hơn; với U-Net (vốn chưa hội tụ), điều này làm khoảng cách lớn thêm.

\subsection{Kiểm chứng hai tuyên bố của paper}
Hai thí nghiệm ở Mục~\ref{sec:robust} dùng 1.000 ảnh test (cứ 3 ảnh lấy 1 theo thứ tự danh sách test, \code{--subset 1000}) để kịp thời gian;
ở $384\times384$ mIoU trên tập con này là 89,0\% (SegFormer) và 43,4\% (U-Net), gần với kết quả trên toàn tập test.

\textbf{Đổi độ phân giải test (Bảng~\ref{tab:resolution}).} Paper cho rằng Mix-FFN thay cho positional encoding giúp mô hình
chịu được độ phân giải test khác lúc train. Trên tập của nhóm, SegFormer giữ mIoU trong khoảng 86,6--89,0\% khi cạnh ảnh đi từ
256 đến 640 (gấp 2,5 lần), tức giảm tối đa 2,4 điểm. Điều này khớp với Bảng 1c của paper (Mix-FFN chỉ giảm 0,7 điểm khi tăng
độ phân giải trên Cityscapes). U-Net thay đổi nhiều hơn theo một hướng đáng chú ý: mIoU \emph{tăng} khi ảnh nhỏ hơn (48,5\% ở 256)
và giảm mạnh khi ảnh lớn hơn (32,3\% ở 640), trong khi IoU thú/nền gần như không đổi (95--97\%). Khi ảnh nhỏ, vùng nhìn của
mỗi điểm ảnh phủ được nhiều phần con vật hơn nên U-Net nhận giống tốt hơn; điều này củng cố giải thích rằng điểm yếu của
U-Net ở đây là thiếu ngữ cảnh toàn cục, không phải thiếu chi tiết.

\textbf{Ảnh hỏng (Bảng~\ref{tab:corruption}, Hình~\ref{fig:robust}).} Paper báo cáo SegFormer-B5 bền hơn hẳn CNN trên
Cityscapes-C. Kết quả của nhóm xác nhận xu hướng này ở cỡ B2 và trên tập khác:
\begin{itemize}
  \item \st{Nhiễu Gauss:} SegFormer gần như không đổi ở $\sigma = 0{,}04$ và $0{,}08$ (88,9\% và 87,5\%), còn 79,5\% ở
        $\sigma = 0{,}16$. U-Net giảm từ 43,4\% xuống 2,6\%: ở mức nặng nhất nó gần như không còn thấy con vật (IoU thú/nền
        5,2\%). Đây là khác biệt lớn nhất, giống với paper (57,8 so với 6,5 trên Cityscapes-C).
  \item \st{Làm mờ và nén JPEG:} cả hai đều giảm mạnh ở mức nặng nhất (SegFormer còn 47,2\% và 45,6\%), nhưng tỉ lệ giữ lại của
        SegFormer vẫn cao hơn: ở mức nặng nhất SegFormer giữ khoảng 53\% và 51\% mIoU ban đầu, U-Net chỉ 28\% và 8\%.
\end{itemize}
Vì U-Net ở đây chưa hội tụ, phép so sánh độ bền giữa hai kiến trúc không hoàn toàn công bằng; kết luận chắc chắn hơn là
\st{SegFormer-B2 tự nó rất bền với nhiễu và đổi độ phân giải}, còn với làm mờ và JPEG nặng thì không mô hình nào bền.

\subsection{Đối chiếu với paper: tái hiện được gì, khác gì, học được gì}
\begin{table}[H]
\centering
\caption{Tuyên bố của paper và kết quả của nhóm.}\label{tab:compare}
\begin{tabularx}{\linewidth}{@{}P{4.4cm}LP{2.2cm}@{}}
\toprule
Tuyên bố của paper & Kết quả của nhóm & Tái hiện? \\\midrule
Encoder MiT + bộ giải mã MLP nhẹ cho kết quả tốt hơn CNN cùng cỡ (Bảng 1d: MiT-B2 45,4 vs ResNet-50 34,7 trên ADE20K) & SegFormer-B2 89,2\% vs U-Net + ResNet-50 46,1\% mIoU; tốt hơn ở cả 37 giống & Có, nhưng khoảng cách bị phóng đại do U-Net chưa hội tụ \\
Không cần positional encoding, chịu được đổi độ phân giải (Bảng 1c) & mIoU giảm tối đa 2,4 điểm khi test từ 256 đến 640 & Có \\
Bền với ảnh hỏng, đặc biệt nhiễu (Bảng 5, B5) & Ở B2: nhiễu nặng còn 79,5\%; làm mờ / JPEG nặng còn 45--47\% & Một phần: bền với nhiễu, không bền với mờ / JPEG nặng \\
Hiệu quả: ít tham số, nhanh & Ít tham số hơn nhưng chậm hơn 1,7 lần và tốn gấp 1,8 lần bộ nhớ trên T4 & Không (trên phần cứng của nhóm) \\
\bottomrule
\end{tabularx}
\end{table}

\textbf{Khác biệt so với paper.} Nhóm huấn luyện khoảng 3.100 bước (paper 160.000), ảnh $384\times384$, một GPU T4, không có
augmentation đổi tỉ lệ 0,5--2,0; bài toán là một con vật trên nền thay vì cảnh nhiều vật. Con số tuyệt đối vì vậy không so
trực tiếp với paper được; nhóm chỉ đối chiếu \emph{xu hướng}.

\textbf{Điều nhóm học được.} (1) Một công thức huấn luyện ``chuẩn'' của một paper có thể bất lợi cho kiến trúc khác: so sánh
công bằng cần tinh chỉnh riêng cho từng mô hình hoặc ít nhất chạy đến khi hội tụ. (2) Tách metric (IoU thú/nền, độ chính
xác giống theo ảnh, nhầm theo loài) giúp chỉ ra chính xác lỗi nằm ở nhận giống chứ không ở khoanh vùng; chỉ nhìn mIoU sẽ không
thấy điều này. (3) ``Ít tham số'' khác ``rẻ'': cần đo thời gian và bộ nhớ thực trên phần cứng định dùng.

% ======================================================================= 8
\section{Hạn chế}
\begin{itemize}
  \item \st{U-Net chưa hội tụ.} Đường mIoU validation của U-Net vẫn tăng ở epoch cuối; giới hạn GPU chỉ cho phép 15 epoch và
        một công thức chung. Thí nghiệm công bằng hơn nên tìm learning rate riêng cho U-Net (ví dụ $10^{-4}$ cho encoder) và
        chạy lâu hơn.
  \item \st{Mỗi cấu hình chỉ chạy một lần} (seed 42), nên không có khoảng tin cậy. Với SegFormer, mIoU validation dao động
        khoảng 1--2 điểm giữa các epoch cuối, nên các chênh lệch nhỏ (ví dụ giữa 100\% và 50\% dữ liệu: 89,2 so với 88,2)
        không chắc có ý nghĩa.
  \item \st{Đánh giá thêm trên 1.000 ảnh test} thay vì toàn bộ 3.662 ảnh do thiếu thời gian GPU.
  \item \st{Nhãn:} vùng viền dày bị bỏ qua nên metric không đo được độ sắc nét của biên; tập test còn ảnh có trimap đảo ngược.
  \item \st{Số epoch giữ cố định cho mọi lượng dữ liệu}, nên cấu hình ít ảnh có ít bước cập nhật hơn.
  \item Chưa làm thử chéo trên Pascal VOC và Grad-CAM như dự định mở rộng ở A2.1.
\end{itemize}

% ======================================================================= 9
\section{Sử dụng công cụ AI}
Theo quy định liêm chính học thuật của môn học, nhóm ghi rõ: trợ lý AI (Claude) được dùng để hỗ trợ viết khung mã nguồn,
soạn bố cục báo cáo và rà soát lỗi. Nhóm đã đọc, chạy kiểm tra, chỉnh sửa và chịu trách nhiệm về toàn bộ nội dung;
mọi lựa chọn chủ đề, paper, dataset và mô hình đều được nhóm hiểu và có thể giải thích.

% ======================================================================= Tài liệu tham khảo
\begin{thebibliography}{9}
\bibitem{xie2021} E. Xie, W. Wang, Z. Yu, A. Anandkumar, J. M. Alvarez, P. Luo. SegFormer: Simple and Efficient Design for Semantic Segmentation with Transformers. \emph{NeurIPS}, 2021. \url{https://arxiv.org/abs/2105.15203}.
\bibitem{parkhi2012} O. M. Parkhi, A. Vedaldi, A. Zisserman, C. V. Jawahar. Cats and Dogs. \emph{CVPR}, 2012.
\bibitem{ronneberger2015} O. Ronneberger, P. Fischer, T. Brox. U-Net: Convolutional Networks for Biomedical Image Segmentation. \emph{MICCAI}, 2015.
\bibitem{he2016} K. He, X. Zhang, S. Ren, J. Sun. Deep Residual Learning for Image Recognition. \emph{CVPR}, 2016.
\bibitem{everingham2010} M. Everingham, L. Van Gool, C. K. I. Williams, J. Winn, A. Zisserman. The Pascal Visual Object Classes (VOC) Challenge. \emph{IJCV}, 2010.
\bibitem{timm} R. Wightman. PyTorch Image Models (timm). \url{https://github.com/huggingface/pytorch-image-models}.
\bibitem{wolf2020} T. Wolf et al. Transformers: State-of-the-Art Natural Language Processing. \emph{EMNLP (System Demonstrations)}, 2020.
\end{thebibliography}

\end{document}
